
\begin{thanks}
首先感谢我的导师马文淦教授的谆谆教诲和悉心指导。在科大的五年学习生涯中，马老师教导我如何选题，
并对我在理论研究和计算中遇到的问题提出了深刻见解和宝贵建议。马老师严谨的治学态度、勤奋的科研习惯，
以及对物理问题的洞察力，都值得我学习，也将继续影响我今后的成长。借此机会，向我最尊敬的马老师表达
诚挚的谢意。今后我将严于律己，努力做出更好的成绩，以此回报老师的教诲。

感谢导师张仁友副教授在工作上给予我的支持与帮助。张老师扎实的理论功底和严谨认真的治学态度，
尤其是在审阅我们的论文时，给我留下了深刻印象。感谢重庆大学的郭磊教授对我们工作中遇到的问题
给予耐心细致的指导。同样感谢蒋一教授和韩良教授对我们工作和生活的关心，感谢班主任李雯老师的热心帮助。

感谢刘文、段鹏飞、李晓周、熊寿健、李霄鹏和陈良文等师兄对我的无私帮助。感谢一起奋斗过的张宇和凌刘生同学
在学习和生活中给予我的帮助。五年来，大家一起学习、一起玩耍、一起成长，谢谢你们！也感谢实验室的陈冲
在我遇到理论问题时给予的帮助。感谢沈永柏、张文娟、何凯、王勇、汪洪、周昊、李合意、杨晓东、周思益和陈畅
等师弟师妹们。

最后，感谢大学同学张敬，他后来也与我一同到科大读研。九年来，我们一起学习，也一起畅谈人生与理想。
感谢大学同学朱春生，他后来在科大代培，在生活和学习中给予我贴心的帮助。感谢我的父母、哥哥和嫂子
对我的理解和帮助，也感谢小侄子给我带来的快乐。

谨以此文献给所有关心、理解、帮助和支持过我的人们！谢谢你们！
\vskip 18pt

\begin{flushright}

~~~~李伟华~~~~

2015年9月25日于中科大近代物理系417实验室

\end{flushright}

\end{thanks}
